\title{DeepSeekMath: Pushing the Limits of Mathematical Reasoning in Open Language Models}

\begin{abstract}
Mathematical reasoning remains a notable obstacle for language models because of its intricate and highly structured demands. This work presents DeepSeekMath~7B, an extension of the DeepSeek-Coder-Base-v1.5 7B model that is further pre-trained on 120B math-centric tokens drawn from Common Crawl, as well as supplementary data in natural language and code. DeepSeekMath~7B attains a noteworthy 51.7\% on the MATH benchmark at the competition level, accomplishing this without recourse to auxiliary toolkits or ensemble (voting) strategies, and reaching a competitive standing with Gemini-Ultra and GPT-4. When utilizing self-consistency sampling over 64 generations, DeepSeekMath~7B raises its score on MATH to 60.9\%. The robust mathematical reasoning proficiency observed in DeepSeekMath~is ascribed to two principal innovations: (1) an elaborate data selection and filtration process that leverages the immense scope of freely available web datasets, and (2) the development of Group Relative Policy Optimization (GRPO), a new approach inspired by Proximal Policy Optimization (PPO), which not only advances reasoning performance in mathematics but also improves PPO's memory efficiency.
\end{abstract}

\section{Introduction}

Large language models (LLMs) have transformed artificial intelligence's capabilities in mathematical reasoning, driving notable progress on benchmarks related to quantitative reasoning \citep{MATH} as well as geometry reasoning \citep{trinh2024solving}. Additionally, these models have demonstrated substantial value in aiding humans with the resolution of challenging mathematical tasks \citep{tao}. Nevertheless, state-of-the-art systems such as GPT-4 \citep{gpt4} and Gemini-Ultra \citep{gemini} remain inaccessible to the public, while open-source alternatives currently available exhibit markedly inferior performance.

This paper presents DeepSeekMath, a specialized language model that markedly surpasses other open-source models in mathematical reasoning, and demonstrates performance on par with GPT-4 on several academic benchmark tasks. Central to this achievement is the construction of the DeepSeekMath Corpus, an expansive, high-quality pre-training dataset consisting of 120B mathematical tokens. The corpus is constructed by leveraging the Common Crawl (CC) and utilizing a classifier based on fastText \citep{joulin2016fasttext}. During the initial phase, the classifier is trained with positive samples taken from OpenWebMath \citep{openwebmath}, complemented by a varied assortment of web pages as negative samples. The classifier is then applied to the CC to extract further relevant data, which is subsequently curated and verified through human annotation. After augmenting the dataset, the classifier undergoes retraining to further enhance its effectiveness. Evaluation outcomes reveal that our corpus maintains exceptional quality, enabling our base model, DeepSeekMath-Base 7B, to reach 64.2\% accuracy on GSM8K \citep{gsm8k} and 36.2\% on the challenging MATH dataset \citep{MATH}, both of which outperform the previous Minerva 540B model \citep{minerva}. The DeepSeekMath Corpus also features multilingual content, which has led to observable gains on Chinese mathematics benchmarks \citep{wei2023cmath,agieval}. We suggest that our strategies for mathematical data curation can inform ongoing research, and there remains considerable potential for subsequent advancements.

DeepSeekMath-Base is built upon DeepSeek-Coder-Base-v1.5 7B \citep{deepseek-coder}, as our findings suggest that initializing from a code-focused model yields better results than using a general-purpose LLM. Additionally, our experiments reveal that mathematical training leads to improvements on the MMLU \citep{mmlu} and BBH \citep{bbh} benchmarks. This suggests that, beyond boosting mathematical proficiency, such training also strengthens the model’s overall reasoning performance.

Following pre-training, we conduct mathematical instruction tuning on DeepSeekMath-Base utilizing datasets incorporating chain-of-thought \citep{cot}, program-of-thought \citep{pot,pal}, and tool-integrated reasoning \citep{tora} approaches. The tuned model, DeepSeekMath-Instruct 7B, surpasses other 7B parameter models and achieves performance on par with open-source instruction-tuned models possessing 70B parameters.

In addition, we present Group Relative Policy Optimization (GRPO), a novel reinforcement learning (RL) algorithm derived from Proximal Policy Optimization (PPO) \citep{schulman2017proximal}. GRPO eliminates the need for a critic by estimating baselines directly from group scores, thereby markedly decreasing the computational resources required for training. Using only a selection of English instruction tuning data, GRPO delivers notable performance enhancements over the already robust DeepSeekMath-Instruct model, as evidenced by improvements on both in-domain benchmarks (GSM8K: 82.9\% $\rightarrow$ 88.2\%, MATH: 46.8\% $\rightarrow$ 51.7\%) and out-of-domain problems (for instance, CMATH: 84.6\% $\rightarrow$ 88.8\%) during the RL process. Furthermore, we propose a unified analytical framework that encompasses various methods—including Rejection Sampling Fine-Tuning (RFT) \citep{yuan2023scaling}, Direct Preference Optimization (DPO) \citep{dpo}, PPO, and our proposed GRPO—demonstrating that these techniques can all be viewed as direct or streamlined approaches to RL. We perform comprehensive experimental analyses, comparing factors such as online versus offline regimes, supervising outcome versus reasoning process, and contrasting single-turn with iterative RL methodologies, to rigorously examine the core aspects of the unified paradigm. Finally, we elucidate how our RL methodology amplifies the effectiveness of instruction-tuned models, and outline promising avenues for further advances in RL within this framework.

\subsection{Contributions}
We present scalable math pre-training and provide an investigation and evaluation of reinforcement learning methodologies.

\textbf{Math Pre-Training at Scale}

\begin{itemize}[topsep=0pt]
    \item This study demonstrates that the openly available Common Crawl dataset includes significant resources relevant to mathematics. Through a carefully engineered data filtering pipeline, we build the DeepSeekMath Corpus—a robust collection featuring 120B tokens extracted from web pages specifically curated for mathematical relevance. This corpus is nearly sevenfold larger than the mathematical web pages utilized in Minerva \citep{minerva}, and approximately nine times greater than those in the newly released OpenWebMath \citep{openwebmath}.

\item The DeepSeekMath-Base 7B pre-trained base model demonstrates performance on par with Minerva 540B \citep{minerva}, suggesting that model size alone does not determine proficiency in mathematical reasoning. High-quality pre-training data enables smaller models to attain competitive results.

\item We present results from our experiments on mathematical training. Conducting code training before engaging in math training enhances models' performance on mathematical problem-solving tasks, irrespective of whether external tools are utilized. This provides some evidence relevant to the ongoing debate: \textit{does code training enhance reasoning skills?} Our findings indicate that it does, particularly in the context of mathematical reasoning.

\item While it is typical to use arXiv papers for training—an approach prevalent in numerous mathematics-focused studies—this strategy does not yield significant gains across any of the mathematical benchmarks utilized in this work.
\end{itemize}

\textbf{Exploration and Analysis of Reinforcement Learning}

\begin{itemize}[topsep=0pt]
    \item We present Group Relative Policy Optimization (GRPO), a reinforcement learning approach designed for both efficiency and efficacy. GRPO eliminates the need for a critic model by deriving the baseline directly from group scores, which notably decreases the computational demands in comparison with Proximal Policy Optimization (PPO).
    \item Our experiments reveal that GRPO offers considerable improvements in the instruction-tuned performance of DeepSeekMath-Instruct when utilizing only instruction-tuning datasets. Additionally, reinforcement learning with GRPO yields better generalization to out-of-domain scenarios.
    \item We establish a comprehensive framework for interpreting and relating various algorithms, including RFT, DPO, PPO, and GRPO. Extensive empirical studies—such as comparisons between online and offline learning, outcome-based versus process-based supervision, and single-step versus iterative reinforcement learning—enable in-depth analysis of the critical components underlying this framework.
    \item Drawing on our unified framework, we analyze the mechanisms contributing to the success of reinforcement learning and outline several promising avenues to further advance reinforcement learning methodologies for LLMs.
\end{itemize}

\subsection{Summary of Evaluations and Metrics}

\begin{itemize}[topsep=0pt]
    \item \textbf{English and Chinese Mathematical Reasoning}: We perform thorough evaluations of our models across both English and Chinese datasets, which encompass mathematical questions ranging from elementary to collegiate complexity. The English datasets used for evaluation are GSM8K \citep{gsm8k}, MATH \citep{MATH}, SAT \citep{llemma}, OCW Courses \citep{minerva}, and MMLU-STEM \citep{mmlu}. For Chinese, we assess performance using MGSM-zh \citep{mgsm}, CMATH \citep{wei2023cmath}, Gaokao-MathCloze \citep{agieval}, and Gaokao-MathQA \citep{agieval}. We examine the models' proficiency in producing independent textual solutions without external computational aids, as well as their effectiveness in resolving tasks through Python programming.

In evaluations on English benchmarks, DeepSeekMath-Base demonstrates performance on par with the proprietary Minerva 540B model \citep{minerva}, while outperforming all open-source base models—including Mistral 7B \citep{mistral} and Llemma-34B \citep{llemma}—by considerable margins, irrespective of whether these models underwent specialized math pre-training. Particularly, DeepSeekMath-Base achieves even better results on Chinese benchmarks. This advantage is likely attributable to our broader approach to data collection: unlike prior studies \citep{minerva, llemma}, which rely exclusively on English-language math data, our pre-training corpus incorporates high-quality datasets from multiple languages. Further enhancements via mathematical instruction tuning and reinforcement learning yield the DeepSeekMath-Instruct and DeepSeekMath-RL models, which deliver state-of-the-art results, exceeding 50\% accuracy on the MATH dataset—a competition-level benchmark—a milestone for open-source models.

\item \textbf{Formal Mathematics}: We assess the performance of DeepSeekMath-Base on the informal-to-formal theorem proving task introduced in \citep{dsp_proof}, utilizing the miniF2F dataset \citep{minif2f} and selecting Isabelle \citep{isabelle} as the proof assistant. The results indicate that DeepSeekMath-Base achieves notable few-shot autoformalization capabilities.
\item \textbf{Natural Language Understanding, Reasoning, and Code}: To thoroughly gauge the models' abilities in general comprehension, reasoning, and code generation, we test DeepSeekMath-Base using the Massive Multitask Language Understanding (MMLU) benchmark \citep{mmlu}, which features 57 multiple-choice assessments across a wide range of disciplines, as well as BIG-Bench Hard (BBH) \citep{bbh}, a suite of 23 tasks focused primarily on complex multi-step reasoning. Additionally, HumanEval \citep{codex} and MBPP \citep{mbpp} are employed to measure code generation proficiency. Pre-training on mathematical content yields improvements in both language comprehension and reasoning tasks.

\section{Math Pre-Training}

\subsection{Data Collection and Decontamination} This section describes the methodology employed to build the DeepSeekMath Corpus utilizing Common Crawl data. Figure \ref{fig:data_collect} illustrates the step-by-step pipeline designed for incrementally assembling an extensive mathematical dataset from Common Crawl. The procedure begins with a seed corpus, which consists of a limited yet high-quality math-related dataset. Importantly, this iterative framework can be adapted for use in various other fields, including coding.

Initially, we select OpenWebMath \citep{openwebmath}, a curated resource comprising high-quality mathematical web content, to serve as our foundational seed dataset. Building on this, we utilize fastText \citep{joulin2016fasttext} to identify additional mathematical web pages with characteristics similar to those in OpenWebMath. For model training, 500,000 randomly sampled entries from the seed corpus serve as positive examples, while an equal number of pages from Common Crawl act as negatives. We leverage an open-source fastText implementation\footnote{\url{https://fasttext.cc}}, setting the vector size to 256, learning rate to 0.1, a maximum word n-gram length of 3, requiring at least 3 word occurrences, and performing 3 epochs of training. To manage the extensive Common Crawl dataset, we apply both URL-based deduplication and near-deduplication methods, yielding a collection of 40B HTML web pages. Employing the trained fastText model, we then extract candidate mathematical pages from this deduplicated pool. Subsequently, to ensure content quality, we sort the retrieved pages by their fastText model scores, retaining only the highest-scoring entries. The quantity of data to keep is empirically determined through preliminary training runs with the top 40B, 80B, 120B, and 160B tokens; for the first iteration, we opt to retain the top 40B tokens.

Following the initial round of data gathering, a substantial portion of mathematical web pages remains uncaptured, primarily due to the fastText model being trained on a relatively homogeneous set of positive examples. To remedy this, we expand the range of mathematical web sources included in the seed corpus, aiming to enhance the fastText model’s effectiveness. Our procedure begins by segmenting the entire Common Crawl dataset into distinct domains, where each domain comprises web pages with the same base URL. We then compute, for each domain, the proportion of web pages successfully collected during the first pass. Domains for which more than 10\% of pages have been collected are designated as math-related (for instance, \url{mathoverflow.net}). 

We proceed to manually review and annotate the specific URLs that correspond to mathematical content within these math-related domains (such as \url{mathoverflow.net/questions}). Web pages connected to these URLs that were not gathered in earlier rounds are subsequently incorporated into the seed set. This method allows us to collect a broader array of positive training examples and, as a result, to retrain fastText to capture a greater number of mathematical documents in subsequent iterations. Ultimately, after four rounds of collection, we compile a total of 35.5M mathematical web pages amounting to 120B tokens. By the fourth iteration, we observe that approximately 98\% of the data was already obtained in the third iteration, leading us to terminate further data collection at that point.

In order to prevent contamination of our benchmark datasets, we adopt the approach described in \cite{deepseek-coder}, where web pages are filtered out if they include content from well-known English mathematical benchmarks such as GSM8K \citep{gsm8k}, MATH \citep{MATH}, as well as Chinese benchmarks like CMATH \citep{wei2023cmath} and AGIEval \citep{agieval}. Our filtering process operates as follows: we exclude any portion of text containing a sequence of 10 consecutive tokens that matches identically with any part of the referenced benchmarks. Additionally, for benchmark texts whose length falls below 10 grams but consists of 3 or more grams, we remove web pages based on exact sequence matching to further mitigate data contamination.

\subsection{Validating the Quality of the DeepSeekMath~Corpus}

We run pre-training experiments to investigate how the DeepSeekMath~Corpus is compared with the recently released math-training corpora:

\begin{itemize}[topsep=0pt]
    \item \textbf{MathPile} \citep{mathpile}: a collection composed from multiple sources, including textbooks, Wikipedia, ProofWiki, CommonCrawl, StackExchange, and arXiv, comprising 8.9B tokens, with arXiv providing over 85\% of the content;
    \item \textbf{OpenWebMath} \citep{openwebmath}: a dataset derived from CommonCrawl, specifically selected for mathematical materials, and amounting to 13.6B tokens;
    \item \textbf{Proof-Pile-2} \citep{llemma}: this resource brings together OpenWebMath, AlgebraicStack (which features 10.3B tokens of math-centric code), and arXiv papers (offering 28.0B tokens). In experiments involving Proof-Pile-2, we adopt the arXiv:Web:Code token proportion of 2:4:1, in line with \cite{llemma}.
\end{itemize}

\subsubsection{Training Setting}
\label{sec:quality-policy}
We fine-tune a general-purpose language model comprising 1.3B parameters—architecturally aligned with the DeepSeek LLM series \citep{deepseek-llm}, referred to here as DeepSeek-LLM 1.3B—on mathematical datasets. Each specific corpus is used to individually train a separate model for a total of 150B tokens. All model training is executed via the resource-efficient HAI-LLM \citep{haillm} training infrastructure. Consistent with the established DeepSeek LLM protocol, we adopt the AdamW optimizer \citep{adamW}, setting $\beta_1=0.9$, $\beta_2=0.95$, and $\mathrm{weight\_decay}=0.1$. The learning rate follows a staged schedule: beginning with a 2,000-step warmup to its maximum, then reducing to 31.6\% at the 80\% progress mark, and further tapering down to 10.0\% by 90\% of training completion. The peak learning rate is capped at 5.3e-4, with training executed using batches of 4 million tokens and a context window of 4K.

\subsubsection{Evaluation Results}

\textbf{The DeepSeekMath~Corpus is of high quality, covers multilingual mathematical content, and is the largest in size.}

\begin{itemize}[topsep=0pt]
    \item \textbf{High-quality}:
    We assess downstream capabilities across 8 mathematics benchmarks utilizing few-shot chain-of-thought prompting \cite{cot}. Results presented in Table \ref{tab:corpora_comparison} establish that the model trained with the DeepSeekMath~Corpus consistently outperforms others. As depicted in Figure \ref{fig:corpora_comparisons}, models trained on DeepSeekMath~Corpus exhibit superior results compared to those trained on Proof-Pile-2 after exposure to 50B tokens (equivalent to one complete epoch of Proof-Pile-2), suggesting that DeepSeekMath~Corpus offers higher overall data quality.
    \item \textbf{Multilingual}:
    DeepSeekMath~Corpus includes multilingual content, with English and Chinese being the most prevalent. According to Table \ref{tab:corpora_comparison}, training with DeepSeekMath~Corpus leads to notable improvements in mathematical reasoning abilities for both English and Chinese language tasks. By contrast, the predominantly English-only existing mathematical corpora show minimal gains or even detrimental effects on Chinese mathematical reasoning performance.
    \item \textbf{Large-scale}:
    DeepSeekMath~Corpus significantly surpasses previous mathematical corpora in scale. Figure \ref{fig:corpora_comparisons} illustrates that DeepSeek-LLM 1.3B, when trained with DeepSeekMath~Corpus, attains faster and more sustained progress, as evidenced by the learning trajectory. By comparison, the smaller baseline corpora are subject to repeated epochs during training, causing their corresponding models' performance to plateau rapidly.
\end{itemize}

\subsection{Training and Evaluating DeepSeekMath-Base 7B}

In this section, we present DeepSeekMath-Base 7B, a foundational model distinguished by its robust mathematical reasoning capabilities. DeepSeekMath-Base 7B builds upon the DeepSeek-Coder-Base-v1.5 7B initialization \citep{deepseek-coder} and is further optimized over 500B training tokens. The dataset for training comprises 56\% content from the DeepSeekMath Corpus, 4\% drawn from AlgebraicStack, 10\% from arXiv, 20\% sourced from Github code repositories, and an additional 10\% of natural language data originating from Common Crawl in English and Chinese. Training largely follows the methodology described in Section \ref{sec:quality-policy}, with two key deviations: the peak learning rate is set to 4.2e-4, and we employ a batch size of 10M tokens.

This study offers an in-depth evaluation of DeepSeekMath-Base 7B’s proficiency in mathematics, with particular emphasis on three areas: generating complete mathematical solutions autonomously, employing computational tools to address mathematical challenges, and executing formal theorem-proving tasks. In addition to the mathematical analysis, we present a broader characterization of the base model’s overall abilities, detailing its competence in natural language comprehension, logical reasoning, and programming.

\paragraph{Mathematical Problem Solving with Step-by-Step Reasoning}

We assess the capability of DeepSeekMath-Base in addressing mathematical problems through few-shot chain-of-thought prompting \citep{cot} over eight distinct benchmarks in both English and Chinese. These evaluation benchmarks include quantitative reasoning datasets (such as GSM8K \citep{gsm8k}, MATH \citep{MATH}, and CMATH \citep{wei2023cmath}) as well as multiple-choice formats (like MMLU-STEM \citep{mmlu} and Gaokao-MathQA \citep{agieval}), spanning a broad spectrum of mathematical domains and ranging in complexity from elementary to university level.

As indicated in Table \ref{tab:base_math_cot}, DeepSeekMath-Base 7B consistently achieves the highest scores across all eight evaluated benchmarks among open-source base models—this comparison includes both the prominent general-purpose model Mistral 7B \citep{mistral} and the more recently introduced Llemma 34B \citep{llemma}, which benefits from specialized math training on Proof-Pile-2 \citep{llemma}. Particularly impressive is its performance on the competition-level MATH dataset, where DeepSeekMath-Base exceeds the results of all open-source base models by more than 10\% in absolute terms. Furthermore, it outperforms the closed-source Minerva 540B \citep{minerva}, a model that is 77 times larger, based on PaLM \citep{palm} and further trained on extensive mathematical corpora.

\paragraph{Mathematical Problem Solving with Tool Use} We assess the capability of models to perform mathematically-oriented reasoning tasks on GSM8K and MATH through few-shot program-of-thought prompting, following the approaches of \citet{pot,pal}. The prompt requires models to generate Python code solutions, where they may employ libraries such as \textit{math} and \textit{sympy} to handle complex calculations. The solution to each problem is determined by evaluating the output produced by running the generated code. Results presented in Table \ref{tab:base_math_pot_proof} demonstrate that DeepSeekMath-Base 7B achieves superior performance compared to the former leading model, Llemma 34B.

\paragraph{Formal Mathematics}

The automation of formal proofs offers significant advantages in verifying the correctness and dependability of mathematical arguments, as well as boosting efficiency—an area that has drawn growing interest in recent years. In this study, we assess DeepSeekMath-Base 7B on the informal-to-formal proof generation task from \citep{dsp_proof}, where the objective is to produce a formal proof using an informal problem description, its corresponding formal version, and an informal proof. Our evaluation uses the miniF2F \citep{minif2f} benchmark, designed for Olympiad-level formal mathematics, where formal proofs are constructed in Isabelle for each task utilizing few-shot prompting. Adhering to the protocol outlined in \cite{dsp_proof}, we employ models to produce proof sketches and apply the standard automated prover Sledgehammer \citep{sledgehammer} to complete any incomplete segments. As reported in Table \ref{tab:base_math_pot_proof}, DeepSeekMath-Base 7B achieves notable results in the automatic formalization of proofs.

\paragraph{Natural Language Understanding, Reasoning, and Code}

We assess natural language understanding using MMLU \citep{mmlu}, evaluate reasoning with BBH \citep{bbh}, and test coding skills on HumanEval \citep{codex} and MBPP \citep{mbpp}. According to the results in Table \ref{tab:understanding_reasoning_code}, DeepSeekMath-Base 7B demonstrates marked improvements over its predecessor, DeepSeek-Coder-Base-v1.5 \citep{deepseek-coder}, in both MMLU and BBH tasks, which underscores the beneficial influence of mathematical pretraining on both language comprehension and reasoning abilities. Furthermore, through the integration of code tokens during continual training, DeepSeekMath-Base 7B successfully preserves the coding performance of DeepSeek-Coder-Base-v1.5 when evaluated on HumanEval and MBPP. Taken together, DeepSeekMath-Base 7B achieves notably higher scores than the broadly used Mistral 7B \citep{mistral} across the evaluated reasoning and coding benchmarks.

\section{Supervised Fine-Tuning}

\subsection{SFT Data Curation}

We develop a mathematical instruction-tuning dataset that encompasses both English and Chinese questions, spanning a variety of mathematical domains and a broad range of difficulty. Each problem in the dataset is accompanied by a solution presented in one of several formats: chain-of-thought (CoT) \citep{cot}, program-of-thought (PoT) \citep{pot,pal}, or tool-integrated reasoning \citep{tora}. In total, the dataset comprises 776K training instances.

\begin{itemize}[topsep=0pt]
    \item \textbf{English mathematical datasets}: We provide tool-integrated annotations for problems from GSM8K and MATH, and incorporate a subset of MathInstruct \citep{MathInstruct} in addition to the training portion of Lila-OOD \citep{lila}, where solutions are formulated using either CoT or PoT approaches. Our English dataset encompasses a wide range of mathematical domains, including but not limited to algebra, probability, number theory, calculus, and geometry.
    \item \textbf{Chinese mathematical datasets}: We assemble a collection of K-12 math problems in Chinese that span 76 distinct sub-topics—such as linear equations—with each problem labeled using both CoT and tool-integrated reasoning formats.
\end{itemize}

\subsection{Training and Evaluating DeepSeekMath-Instruct 7B}

This section presents DeepSeekMath-Instruct 7B, an extension of DeepSeekMath-Base that is fine-tuned through mathematical instruction data. To prepare training data, examples are concatenated in a random order up to a context length limit of 4K tokens. The model undergoes 500 training steps, utilizing a batch size of 256 and maintaining a fixed learning rate of 5e-5 throughout.

We evaluate models' mathematical performance both without and with tool use, on 4 quantitative reasoning benchmarks in English and Chinese. We benchmark our model against the leading models of the time:

\begin{itemize}[topsep=0pt]
    \item \textbf{Closed-source models} comprise: (1) the GPT series, with GPT-4 \citep{gpt4} and GPT-4 Code Interpreter~\footnote{\url{https://openai.com/blog/chatgpt-plugins\#\#code-interpreter}} standing out as the leading variants, (2) Gemini Ultra and Gemini Pro \citep{gemini}, (3) Inflection-2 \citep{inflection-2}, (4) Grok-1~\footnote{\url{https://x.ai/model-card}}, along with several models from Chinese firms such as (5) Baichuan-3~\footnote{\url{https://www.baichuan-ai.com}}, and (6) the newest iteration, GLM-4~\footnote{\url{https://open.bigmodel.cn/dev/api\#glm-4}}.

    from the GLM family \citep{glm}.

These models are intended for broad applications and have typically been subjected to various alignment steps.
    \item \textbf{Open-source models} encompass: general-purpose models such as (1) DeepSeek-LLM-Chat 67B \citep{deepseek-llm}, (2) Qwen 72B \citep{qwen}, (3) SeaLLM-v2 7B \citep{seallm}, and (4) ChatGLM3 6B \citep{chatglm3}; in addition, several are tailored for improved mathematical capabilities, including (5) InternLM2-Math 20B~\footnote{\url{https://github.com/InternLM/InternLM-Math}}, an InternLM2 variant trained further on mathematics and subsequently instruction-tuned, (6) Math-Shepherd-Mistral 7B, which uses PPO training \citep{schulman2017proximal} applied to Mistral 7B \citep{mistral} with a process-supervised reward model, (7) the WizardMath line of models \citep{wizardmath}, which augments Mistral 7B and Llama-2 70B \citep{llama2} with evolved-instruct methods (instruction tuning with AI-generated instructions) as well as PPO, predominantly using GSM8K and MATH datasets for training, (8) MetaMath 70B \citep{metamath}, a Llama-2 70B model refined with an expanded GSM8K and MATH dataset, (9) ToRA 34B \cite{tora}, based on CodeLlama 34B, enhanced for mathematical reasoning with integrated tools, and (10) MAmmoTH 70B \citep{MathInstruct}, created by instruction-tuning Llama-2 70B on MathInstruct.

Table \ref{tab:sft_rl_math} illustrates that, when evaluated without access to external tools, DeepSeekMath-Instruct 7B excels in step-by-step mathematical reasoning. Specifically, on the competition-level MATH benchmark, this model exceeds the performance of all other open-source alternatives and outperforms most proprietary systems (including Inflection-2 and Gemini Pro) by at least 9\% in absolute terms. This superior result holds even in comparison to much larger models (such as Qwen 72B) and those refined through reinforcement learning targeted at mathematical tasks (e.g., WizardMath-v1.1 7B). Although DeepSeekMath-Instruct matches the capabilities of major Chinese proprietary models like GLM-4 and Baichuan-3 on the MATH dataset, it still falls short of state-of-the-art proprietary offerings like GPT-4 and Gemini Ultra.

When assessed under conditions permitting models to combine both natural language reasoning and the use of tools via programs, DeepSeekMath-Instruct 7B attains nearly 60\% accuracy on the MATH dataset, outperforming all currently available open-source alternatives. Across additional evaluation benchmarks, our model demonstrates performance on par with DeepSeek-LLM-Chat 67B, the previous leading model that possesses an order of magnitude more parameters.

\subsection{Group Relative Policy Optimization} Reinforcement learning (RL) has demonstrated significant success in enhancing the mathematical reasoning skills of LLMs following the Supervised Fine-Tuning (SFT) phase \citep{wang2023math,wizardmath}. In this part, we present Group Relative Policy Optimization (GRPO), a novel RL method designed for both efficiency and efficacy.

\subsubsection{From PPO to GRPO} Proximal Policy Optimization (PPO) \citep{schulman2017proximal} is an actor-critic RL algorithm that is widely used in the RL fine-tuning stage of LLMs \citep{ouyang2022training}. In particular, it optimizes LLMs by maximizing the following surrogate objective:

\begin{equation}
\footnotesize
    \mathcal{J}_{PPO}(\theta) = \mathbb{E}{[q \sim P(Q), o \sim \pi_{\theta_{old}}(O|q)]} \frac{1}{|o|} \sum_{t=1}^{|o|} \min \left[ \frac{\pi_\theta(o_{t} | q, o_{<t})}{\pi_{\theta_{old}}(o_{t} | q, o_{<t})} A_{t}, \text{clip} \left( \frac{\pi_\theta(o_{t} | q, o_{<t})}{\pi_{\theta_{old}}(o_{t} | q, o_{<t})}, 1 - \varepsilon, 1 + \varepsilon \right)  A_{t} \right] ,
\end{equation}

The above equation defines the PPO objective $\mathcal{J}_{PPO}(\theta)$ as an expectation, where samples $(q, o)$ are drawn such that $q$ comes from the distribution $P(Q)$ and $o$ is sampled from the old policy $\pi_{\theta_{old}}(O|q)$. The per-token contributions are averaged over the sequence length $|o|$. For each time step $t$, the objective evaluates the minimum between the ratio of policy probabilities $\frac{\pi_\theta(o_{t} | q, o_{<t})}{\pi_{\theta_{old}}(o_{t} | q, o_{<t})}$ multiplied by the advantage $A_t$, and the advantage weighted by a clipped policy ratio, $\text{clip}(\cdot, 1 - \varepsilon, 1 + \varepsilon)$, ensuring conservative updates and preventing large policy changes.

In this context, $\pi_{\theta}$ denotes the current policy while $\pi_{\theta_{old}}$ refers to the prior version. The variables $q$ and $o$ represent the questions and generated outputs, respectively, sampled from the question dataset and from the outputs produced by $\pi_{\theta_{old}}$. The hyper-parameter $\varepsilon$ is associated with the clipping mechanism employed in PPO, which is designed to enhance the stability of the training process. The advantage estimate $A_t$ is determined through Generalized Advantage Estimation (GAE) \citep{gae}, utilizing the rewards $\{r_{\ge t}\}$ in conjunction with a learned value function $V_{\psi}$. Consequently, within PPO, the learning of a value function must occur concurrently with the policy model. To reduce the risk of excessive reward model optimization, a standard practice involves incorporating a KL penalty with respect to a reference model on a per-token basis into the reward signal at each step \citep{ouyang2022training}, such that

\begin{equation}
    r_{t} = r_\varphi(q, o_{\le t}) - \beta \log\frac{\pi_{\theta}(o_{t}|q, o_{<t})}{\pi_{ref}(o_{t}|q, o_{<t})},
\label{eq:PPO-reward}
\end{equation}

Here, $r_\varphi$ represents the reward model, $\pi_{ref}$ denotes the reference model—commonly corresponding to the initial SFT model—and $\beta$ signifies the parameter controlling the strength of the KL divergence penalty.

As the value function employed in PPO is typically another model of comparable size as the policy model, it brings a substantial memory and computational burden. Additionally, during RL training, the value function is treated as a baseline in the calculation of the advantage for variance reduction. While in the LLM context, usually only the last token is assigned a reward score by the reward model, which may complicate the training of a value function that is accurate at each token. To address this, as shown in Figure \ref{fig:grpo}, we propose Group Relative Policy Optimization (GRPO), which obviates the need for additional value function approximation as in PPO, and instead uses the average reward of multiple sampled outputs, produced in response to the same question, as the baseline. More specifically, for each question $q$, GRPO samples a group of outputs $\{o_1, o_2, \cdots, o_G\}$  from the old policy  $\pi_{\theta_{old}}$  and then optimizes the policy model by maximizing the following objective:

\begin{equation}
\footnotesize
\begin{split}
    \mathcal{J}_{GRPO}(\theta) &= \mathbb{E}_{q \sim P(Q), \{o_i\}_{i=1}^G \sim \pi_{\theta_{old}}(O|q)}  \\
    & \frac{1}{G}\sum_{i=1}^G\frac{1}{|o_i|} \sum_{t=1}^{|o_i|} \left\{ \min \left[ \frac{\pi_\theta(o_{i,t} | q, o_{i,<t})}{\pi_{\theta_{old}}(o_{i,t} | q, o_{i,<t})} \hat{A}_{i,t},\ \text{clip} \left( \frac{\pi_\theta(o_{i,t} | q, o_{i,<t})}{\pi_{\theta_{old}}(o_{i,t} | q, o_{i,<t})}, 1 - \varepsilon, 1 + \varepsilon \right) \hat{A}_{i,t} \right] - \beta \mathbb{D}_{KL}\left[\pi_{\theta} || \pi_{ref}\right]\right\},
\end{split}
\label{eq:GRPO-obj}
\end{equation}

where $\varepsilon$ and $\beta$ are hyper-parameters, and $\hat{A}_{i,t}$  is the advantage calculated based on relative rewards of the outputs inside each group only, which will be detailed in the following subsections. The group relative way that GRPO leverages to calculate the advantages, aligns well with the comparative nature of rewards models,  as reward models are typically trained on datasets of comparisons between outputs on the same question. Also note that, instead of adding KL penalty in the reward, GRPO regularizes by directly adding the KL divergence between the trained policy and the reference policy to the loss, avoiding complicating the calculation of  $\hat{A}_{i,t}$. And different from the KL penalty term used in (\ref{eq:PPO-reward}), we estimate the KL divergence with the following unbiased estimator \citep{kl_approx}:

\begin{equation}
\small
    \mathbb{D}_{KL}\left[\pi_{\theta} || \pi_{ref}\right] = \frac{\pi_{ref}(o_{i,t}|q,o_{i,<t})}{\pi_{\theta}(o_{i,t}|q,o_{i,<t})} - \log\left(\frac{\pi_{ref}(o_{i,t}|q,o_{i,<t})}{\pi_{\theta}(o_{i,t}|q,o_{i,<t})}\right) - 1,
\end{equation}

which is guaranteed to be positive.

\begin{algorithm}[t]
  \small
  \caption{Iterative Group Relative Policy Optimization}
  \textbf{Input} initial policy model $\pi_{\theta_{\text{init}}}$; reward models $r_\varphi$; task prompts $\mathcal{D}$; 
  hyperparameters $\varepsilon$, $\beta$, $\mu$
  \begin{algorithmic}[1]
    \State Set the current policy model $\pi_\theta$ to $\pi_{\theta_{\text{init}}}$
    \For{iteration = 1, \dots, I}
      \State Assign the reference model $\pi_{ref}$ to the latest $\pi_\theta$
      \For{step = 1, \dots, M}
        \State Draw a mini-batch $\mathcal{D}_b$ from $\mathcal{D}$
        \State Set $\pi_{\theta_{old}}$ equal to the current $\pi_\theta$
        \State For every question $q$ in $\mathcal{D}_b$, generate $G$ outputs $\{o_i\}_{i=1}^G$ using $\pi_{\theta_{old}} (\cdot \mid q)$
        \State Using the reward model $r_{\varphi}$, calculate rewards $\{r_i\}_{i=1}^{G}$ for all generated $o_i$
        \State Determine the group-based relative advantages $\hat{A}_{i,t}$ at token $t$ for each $o_i$
        \For{each GRPO iteration = 1, \dots, $\mu$}
          \State Optimize the policy $\pi_\theta$ with respect to the GRPO objective function (see Equation \ref{eq:GC-GRPO})
        \EndFor
      \EndFor 
      \State Refine $r_\varphi$ further via ongoing training incorporating a replay buffer mechanism.
    \EndFor
  \end{algorithmic}
  \textbf{Output} $\pi_\theta$
  \label{alg:iter-grpo}
\end{algorithm}

\subsubsection{Outcome Supervision RL with GRPO} Specifically, for a given question $q$, a collection of outputs $\{o_1, o_2, \cdots, o_G\}$ is generated by sampling from the prior policy model $\pi_{\theta_{old}}$. Each output is then evaluated by a reward model, producing a set of $G$ reward values $\mathbf{r} = \{r_1, r_2, \cdots, r_G\}$ accordingly. To standardize the rewards, each value is mean-centered and scaled by the standard deviation computed over the group. In outcome supervision, the resulting normalized reward is provided only at the end of each sequence $o_i$ and is assigned as the advantage $\hat{A}_{i, t}$ for all tokens in that sequence, that is, $\hat{A}_{i, t} = \widetilde{r}_i = \frac{r_i- {\rm mean}(\mathbf{r})}{{\rm std}(\mathbf{r})}$. The policy is then updated by optimizing the objective described in equation (\ref{eq:GRPO-obj}).

\subsubsection{Process Supervision RL with GRPO}  
While outcome supervision limits feedback to a single reward at the conclusion of each output, this approach may be inadequate and inefficient for guiding policies tasked with solving intricate mathematical problems. Inspired by \cite{wang2023math}, we investigate the use of process supervision, which delivers feedback at each reasoning stage rather than solely at the final step. Specifically, for a given question $q$ and $G$ generated outputs $\{o_1, o_2, \cdots, o_G\}$, a process reward model evaluates every intermediate step, assigning a set of rewards: $\mathbf{R} = \{ \{r_1^{index(1)},\cdots,r_1^{index(K_1)}\}, \cdots,  \{r_G^{index(1)},\cdots,r_G^{index(K_G)}\} \}$, with $index(j)$ indicating the token at the conclusion of the $j$-th step, and $K_i$ representing the total number of reasoning steps in the $i$-th output. To enable robust comparison, these rewards are standardized according to their mean and standard deviation, so that $\widetilde{r}_i^{index(j)} = \frac{r_i^{index(j)} - {\rm mean(\mathbf{R})}}{{\rm std(\mathbf{R})}}$.  
Next, within the process supervision scheme, the advantage attributed to each token is computed as the cumulative sum of all subsequent normalized step rewards, namely, $\hat{A}_{i, t} = \sum_{index(j) \ge t} \widetilde{r}_i^{index(j)}$.
Finally, policy optimization is conducted by maximizing the objective function delineated in equation (\ref{eq:GRPO-obj}).

\subsubsection{Iterative RL with GRPO} During reinforcement learning, the previous reward model can become inadequate for effectively guiding the training of the current policy model. To address this, we investigate the iterative application of GRPO. As depicted in Algorithm \ref{alg:iter-grpo}, the iterative GRPO approach involves generating updated training datasets for the reward model using samples drawn from the ongoing policy model. The reward model is further updated through continuous training that employs a replay buffer consisting of 10\% previously collected data. Following this, the policy model is assigned as the new reference model, and training of the policy model resumes, now utilizing the updated reward model.

\subsection{Training and Evaluating DeepSeekMath-RL}

We perform reinforcement learning (RL) using the DeepSeekMath-Instruct 7B model. For RL training, we utilize chain-of-thought-style questions from the GSM8K and MATH datasets present in the SFT data, totaling approximately 144,000 questions. Other SFT-derived questions are omitted to specifically assess how RL affects benchmarks that do not have exposure to related data during this phase. The reward model training set is built according to the methodology outlined in \citep{wang2023math}. Our initial reward model is trained with DeepSeekMath-Base 7B, using a learning rate of $2\times10^{-5}$. In the GRPO process, we set the policy model's learning rate to $1\times10^{-6}$, and use a KL regularization coefficient of $0.04$. For each question, $64$ samples are generated. The sequence length and the training batch size are both set to 1024. After every exploration stage, the policy model undergoes just one update. Evaluation of DeepSeekMath-RL 7B is performed using the same benchmarks as for DeepSeekMath-Instruct 7B. In this setup, the GSM8K and MATH datasets—with their chain-of-thought reasoning—are treated as in-domain, while the rest of the benchmarks serve as out-of-domain evaluations.

As shown in Table \ref{tab:sft_rl_math}, both open- and closed-source models were evaluated using chain-of-thought and tool-integrated reasoning approaches on English and Chinese datasets. Our analysis reveals the following: (1) With chain-of-thought reasoning, DeepSeekMath-RL 7B achieves accuracy rates of 88.2\% on GSM8K and 51.7\% on MATH, exceeding the results of all other open-source models within the 7B to 70B parameter range and outperforming most closed-source competitors as well. (2) Notably, the training of DeepSeekMath-RL 7B was limited solely to chain-of-thought-format instruction tuning data for GSM8K and MATH, building upon DeepSeekMath-Instruct 7B. Even with this restricted training corpus, DeepSeekMath-RL 7B consistently surpasses DeepSeekMath-Instruct 7B on every evaluation metric, highlighting the advantages conferred by reinforcement learning.

\section{Discussion}
This section presents an overview of the outcomes from both the pre-training phase and the reinforcement learning (RL) experiments.

\subsection{Lessons Learnt in Pre-Training}

We begin by describing our pre-training process. Except where indicated, we follow the training configurations described in Section \ref{sec:quality-policy}. It should be emphasized that throughout this section, references to the DeepSeekMath Corpus pertain to an 89B-token dataset compiled during the second round of data acquisition.

\subsubsection{Code Training Benefits Mathematical Reasoning}

An often-cited but unsubstantiated hypothesis claims that code training enhances reasoning skills. In this work, we address this claim in part, focusing specifically on mathematics: code training enhances models’ mathematical reasoning performance, regardless of whether external tools are involved.

To study how code training affects mathematical reasoning, we experimented with the following two-stage training and one-stage training settings:

\textbf{Two-Stage Training}

\begin{itemize}[topsep=0pt]
    \item \textbf{Code Training for 400B Tokens $\rightarrow$ Math Training for 150B Tokens}: Our approach involves first pre-training DeepSeek-LLM 1.3B on 400B tokens sourced from code datasets, then continuing with an additional 150B tokens drawn from math-related data.
    \item \textbf{General Training for 400B Tokens $\rightarrow$ Math Training for 150B Tokens}: To provide a baseline for comparison, we conduct an experiment where the model is initially exposed to 400B tokens from a comprehensive general text corpus (curated by DeepSeek-AI), rather than code, before proceeding to the same 150B token math-focused training, aiming to assess the relative impact of code versus general pre-training on enhancing mathematical reasoning capabilities.
\end{itemize}

\textbf{One-Stage Training}

\begin{itemize}[topsep=0pt]
    \item \textbf{Math Training on 150B Tokens}: DeepSeek-LLM 1.3B undergoes training on a dataset consisting of 150B math-specific tokens;
    \item \textbf{Joint Training with 400B Code Tokens and 150B Math Tokens}: Sequentially training on math tokens after code tokens results in diminished code capabilities. To address this, we explore the impact of simultaneously introducing both code and math tokens in a unified training phase, aiming to enhance mathematical reasoning while mitigating catastrophic forgetting.
\end{itemize}

\paragraph{Results}
The downstream outcomes across various training configurations are displayed in Table \ref{tab:code-math} and Table \ref{tab:code-math-general-eval}.

Program-aided mathematical reasoning benefits from code training, regardless of whether it is implemented with two-stage or one-stage training procedures. As indicated in Table \ref{tab:code-math}, code training by itself in the two-stage framework substantially improves performance on GSM8K and MATH tasks that utilize Python. Adding math training as a subsequent stage leads to additional performance gains. Notably, in the one-stage training configuration, jointly incorporating code tokens and math tokens successfully alleviates the problem of catastrophic forgetting often encountered with two-stage training. Furthermore, this approach enhances both coding performance (as seen in Table \ref{tab:code-math-general-eval}) and program-aided mathematical reasoning capabilities (Table \ref{tab:code-math}).

Training with code data also enhances mathematical reasoning skills even when external tools are not employed. In the two-stage training framework, the first phase that involves code learning alone already yields noticeable improvements. Furthermore, this initial code-focused phase increases the effectiveness of later math-specific training, ultimately achieving optimal results. In contrast, merging code tokens with math tokens in a single-stage training approach leads to diminished mathematical reasoning abilities in the absence of tool assistance. A possible explanation is that DeepSeek-LLM 1.3B may not have sufficient scale or capacity to thoroughly learn from both code and math datasets at the same time.

\subsubsection{ArXiv Papers Seem Ineffective in Improving Mathematical Reasoning}

ArXiv papers are commonly included as a component of math pre-training data \citep{minerva,gpt-f,llemma,mathpile}. However, detailed analysis regarding their impact on mathematical reasoning has not been extensively conducted. Perhaps counter-intuitively, according to our experiments, arXiv papers seem ineffective in improving mathematical reasoning. We experiment with models of different sizes, including DeepSeek-LLM 1.3B and DeepSeek-Coder-Base-v1.5 7B \citep{deepseek-coder}, using arXiv corpora that underwent varied processing pipelines:

\begin{itemize}[topsep=0pt]
    \item \textbf{MathPile} \citep{mathpile}: a dataset comprising 8.9 billion tokens, assembled using heuristic-based cleaning and filtering methods, and primarily sourced (over 85\%) from scholarly arXiv publications;
    \item \textbf{ArXiv-RedPajama} \citep{redpajama}: a complete collection of arXiv LaTeX documents that have had preambles, comments, macros, and bibliography sections excluded, amounting to a total of 28.0 billion tokens.
\end{itemize}

In our study, we conducted separate training sessions where DeepSeek-LLM 1.3B processed 150B tokens and DeepSeek-Coder-Base-v1.5 7B was exposed to 40B tokens, each utilizing distinct arXiv corpora. Our findings suggest that incorporating arXiv papers does not enhance a model’s capacity for mathematical reasoning. Training exclusively on arXiv data yielded no significant gains—or even declines—in model performance across a range of mathematical benchmarks with varying difficulty levels considered in this work. These assessments comprised quantitative reasoning datasets such as GSM8K and MATH (see Table \ref{tab:arxiv-cot}), multiple-choice tasks like MMLU-STEM (Table \ref{tab:arxiv-cot}), and formal mathematics evaluations including miniF2F (refer to Table \ref{tab:arxiv_atp}).

However, this conclusion has its limitations and should be taken with a grain of salt. We have not yet studied:

\begin{itemize}[topsep=0pt]
    \item The influence of arXiv tokens on particular mathematics-focused applications outside the scope of this study, for example, transforming formal theorems or proofs into their informal counterparts (informalization);
    \item The role arXiv tokens play when integrated with additional data sources;
    \item If the advantages provided by arXiv documents would be observable with models of greater scale.
\end{itemize}

Thus, further exploration is required, which we leave for future studies.

\subsection{Insights of Reinforcement Learning}
\subsubsection{Towards to a Unified Paradigm} This section introduces a comprehensive framework designed to examine various training strategies, including SFT, RFT, DPO, PPO, and GRPO. Additionally, we undertake experiments to investigate key components within this unified framework. Typically, the gradient with respect to the parameter $\theta$ for a given training approach is formulated as follows:

\begin{equation}
    \nabla_{\theta}\mathcal{J}_{\textcolor{red}{\mathcal{A}}}(\theta) = \mathbb{E}[\underbrace{(q,o) \sim \textcolor{red}{\mathcal{D}}}_{Data \ Source}]\left( \frac{1}{|o|} \sum_{t=1}^{|o|}  \underbrace{GC_{{\mathcal{A}}}(q, o, t, \textcolor{red}{\pi_{{rf}}})}_{Gradient \ Coefficient}  \nabla_{\theta}\log \pi_{\theta}(o_t | q, o_{<t})\right).
\label{eq:objective}
\end{equation}

There exist three key components: 1) \textit{Data Source $\mathcal{D}$}, which determines the training data; 2) \textit{Reward Function $\pi_{{rf}}$}, which is the source of the training reward signal; 3) \textit{Algorithm $\mathcal{A}$}: which processes the training data and the reward signal to the gradient coefficient $GC$ that determines the magnitude of the penalty or reinforcement for the data. We analyze several representative methods based on such a unified paradigm:

\begin{itemize}
    \item  \textbf{Supervised Fine-tuning (SFT)}: SFT involves adjusting a pretrained model using a dataset curated by humans specifically for fine-tuning purposes.
    \item \textbf{Rejection Sampling Fine-tuning (RFT)}: RFT extends SFT by further training the SFT model on selected responses, which are generated by the SFT model in response to SFT questions. These responses are filtered to retain only those with correct answers.
    \item \textbf{Direct Preference Optimization (DPO)}: DPO enhances the SFT model through additional fine-tuning, utilizing paired outputs generated by the SFT model and optimizing them according to a pair-wise DPO loss function.
    \item \textbf{Online Rejection Sampling Fine-tuning (Online RFT)}: Unlike standard RFT, Online RFT starts with the SFT model as the policy model, and continues refinement by fine-tuning using augmented samples generated from the currently updated policy model.
    \item \textbf{PPO/GRPO}: PPO/GRPO methods begin by initializing the policy model with the SFT model and proceed by further reinforcing the model through samples produced by the evolving policy model in real time.
\end{itemize}

The elements comprising these approaches are outlined in Table \ref{tab:data_policy}. For an in-depth discussion of the derivation procedure, consult Appendix \ref{app:analysis-rl}.

\paragraph{Observation about Data Source} We classify data sources into two types: online sampling and offline sampling. Online sampling refers to collecting training data generated by the policy model as it interacts with the environment in real time during training. In contrast, offline sampling involves using training data obtained from the outputs of the pre-trained SFT model. Both RFT and DPO utilize the offline sampling approach, whereas Online RFT and GRPO employ online sampling methods.

As illustrated in Figure \ref{fig:rl-analysis}, our results indicate that Online RFT achieves markedly better performance than RFT across both benchmarks. In particular, the performance of Online RFT matches that of RFT during the early phases of training, but then surpasses RFT significantly as training progresses. This outcome highlights the benefits of an online training approach. This finding aligns with expectations: early in training, the actor remains quite similar to the SFT model, and consequently, samples generated reflect only slight distinctions. As training advances, however, samples from the actor diverge more substantially, making real-time data collection increasingly beneficial.

\paragraph{Observation about Gradient Coefficient} The process by which the algorithm incorporates the reward signal into the gradient coefficient is central to updating the model parameters. In our experimental setup, we differentiate the reward function into two types: `Rule' and `Model'. The `Rule' method involves assessing response quality through evaluation of answer correctness, whereas the `Model' approach involves developing a reward model to assign scores to responses, where the reward model itself is trained on data labeled using the rule-based method. Equations \ref{eq:GC-RFT} and \ref{eq:GC-GRPO} illustrate a fundamental distinction between GRPO and Online RFT. Specifically, GRPO individually calibrates its gradient coefficient in response to the reward value assigned by the reward model, thereby providing varying degrees of reinforcement or penalization depending on the magnitude of the response scores. In sharp contrast, Online RFT does not possess this property; it fails to penalize incorrect responses and instead applies a uniform reinforcement strength to all responses deemed correct.

As shown in Figure \ref{fig:rl-analysis}, GRPO outperforms online RFT, emphasizing the advantages of modifying the coefficients associated with positive and negative gradients. Moreover, GRPO+PS achieves better results than GRPO+OS, which suggests that leveraging more granular, step-sensitive gradient coefficients is beneficial. Additionally, we investigate iterative RL by conducting two sequential iterations in our experiments. Figure \ref{fig:iter-rl} reveals that iterative RL leads to substantial performance gains, with the most pronounced improvement occurring during the initial iteration.

\subsubsection{Why RL Works?} In this study, reinforcement learning is applied using a portion of the instruction tuning dataset, resulting in marked performance gains over the baseline instruction-tuned model. To elucidate the mechanisms underlying the effectiveness of reinforcement learning, we assess the Pass@K and Maj@K accuracies of both the Instruct and RL-optimized models across two standard benchmarks. As presented in Figure \ref{fig:maj-pass}, reinforcement learning notably raises Maj@K accuracy while leaving Pass@K largely unchanged. This suggests that RL primarily strengthens the reliability of the model's output distribution, implying that \textbf{the observed improvement may arise from a higher likelihood of generating a correct response among the TopK candidates, rather than from fundamental enhancements in model ability.} Analogously, \citep{wang2023making} described a \textbf{misalignment problem} observed in SFT models on reasoning tasks, and demonstrated that strategies focusing on preference alignment \citep{yuan2023rrhf,song2023preference,wang2023making} can lead to measurable gains in SFT model reasoning performance.

\subsubsection{How to Achieve More Effective RL?}
\label{sec:effec_rl}
Our experiments show that RL performs effectively in the context of mathematical reasoning tasks. Additionally, we introduce a unified framework that allows various prominent training approaches to be analyzed within a common perspective. In this framework, all methods can be interpreted as either straightforward or simplified applications of RL. As outlined in Equation \ref{eq:objective}, three principal elements are identified: Data Source, Algorithm, and Reward Function. We outline several prospective research directions pertaining to these three factors.

\paragraph{Data Source} The data source serves as the foundational input for all training approaches. Within the RL framework, we define the data source as consisting of unlabeled questions alongside outputs generated through sampling from the policy model. For this study, our experiments utilize only the instruction tuning stage's questions, and outputs are sampled using a basic nucleus sampling approach. We hypothesize that this limitation may contribute to our RL pipeline's performance gains being confined to Maj@K metrics. Looking forward, we aim to extend our investigation of the RL pipeline to address out-of-distribution prompts, and we intend to integrate \textbf{advanced sampling (decoding) strategies} such as those employing tree-search algorithms \citep{yao2023tree}. Moreover, \textbf{efficient inference techniques} \citep{xia-etal-2023-speculative,leviathan2023fast,kwon2023efficient,xia2024unlocking}, which significantly influence how effectively the policy model explores, are anticipated to have a major impact as well.

\paragraph{Algorithms} Algorithms utilize data alongside the reward signal in order to determine the gradient coefficients used for model parameter updates. Referring to Equation \ref{eq:objective}, existing approaches generally exhibit a strong tendency to \textbf{TRUST} the reward function’s guidance, either promoting or discouraging the conditional likelihood of particular tokens. Nonetheless, this trust is problematic as the reward signal's accuracy cannot always be assured, particularly in highly intricate tasks. For instance, the PRM800K datasets \citep{lightman2023let}, despite being painstakingly labeled by expert annotators, still feature an error rate of approximately 20\% in their annotations\footnote{\url{https://github.com/openai/prm800k/issues/12\#issuecomment-1728491852}}. Motivated by such limitations, we seek to investigate reinforcement learning algorithms that maintain robustness when faced with unreliable, noisy rewards. We anticipate that adopting \textbf{WEAK-TO-STRONG} \citep{burns2023weak} alignment strategies could profoundly reshape the foundations of learning algorithm design.

\paragraph{Reward Function} The reward function delivers the fundamental training feedback in RL. Typically, this role is filled by a neural reward model. We identify three principal research avenues for reward models: 1) \textbf{Strategies to improve the reward model's generalization capability.} It is imperative that the reward model is robust when confronted with out-of-distribution prompts and complex decoding scenarios; if this condition fails, reinforcement learning may only regularize the output distributions of LLMs without contributing to meaningful capability gains; 2) \textbf{Approaches for representing reward model uncertainty.} Incorporating measures of uncertainty could serve to connect limited reward models with algorithms that learn from weak to strong supervision; 3) \textbf{Methods for constructing high-quality process reward models efficiently}, enabling the provision of detailed feedback throughout the reasoning trajectory \citep{lightman2023let,wang2023math}.

\section{Conclusion, Limitation, and Future Work}

This work introduces DeepSeekMath, which achieves state-of-the-art results among open-source models on the competition-level MATH benchmark and rivals the accuracy of proprietary alternatives. DeepSeekMath~leverages DeepSeek-Coder-v1.5 7B as its foundation, receiving an additional 500B tokens of continual training, notably including 120B math-related tokens extracted from Common Crawl. Through comprehensive ablation experiments, we find that web-sourced pages represent a substantial resource for high-quality mathematical datasets, whereas the benefits derived from arXiv data are less pronounced than initially anticipated. We propose Group Relative Policy Optimization (GRPO), an adaptation of Proximal Policy Optimization (PPO), demonstrating its ability to enhance mathematical reasoning while requiring reduced memory resources. Experimental findings confirm that GRPO yields significant gains even when DeepSeekMath-Instruct 7B attains elevated benchmark results. Additionally, we provide a cohesive framework for conceptualizing a family of related reinforcement learning techniques and identify several promising future research avenues for improved algorithmic performance.

While DeepSeekMath~demonstrates strong results on benchmarks focused on quantitative reasoning, its performance in geometry and theorem-proving falls short compared to closed models. In our preliminary assessment, the model struggled with tasks involving triangles and ellipses, which might reflect biases in both pre-training and fine-tuning data. Furthermore, due to limitations in model size, DeepSeekMath~lags behind GPT-4 with respect to few-shot learning abilities. Whereas GPT-4 benefits noticeably from few-shot prompts, DeepSeekMath~yields nearly the same results in both zero-shot and few-shot settings. Going forward, we intend to refine our data selection processes to build a more robust and diverse pre-training dataset. Moreover, we plan to investigate the avenues (Section \ref{sec:effec_rl}) for enhancing reinforcement learning methods for LLMs.

\end{document}